\section{Introduction}

The continued deployment of low-Earth-orbit (LEO) satellite constellations is increasing the capacity required from feeder links between satellites and optical ground stations~\cite{lebidan2023frame}. Coherent free-space optical (FSO) communication is attractive for this role because coherent reception offers high sensitivity, supports spectrally efficient higher-order modulation, and enables digital compensation of carrier impairments~\cite{valjus2025dsp,paillier2020spaceground,guiomar2022coherent}. The atmospheric path nevertheless introduces absorption, scattering, and refractive-index fluctuations. Turbulence produces scintillation in the received irradiance and perturbs the optical phase, while satellite motion and finite laser linewidth add time-varying carrier-frequency and phase errors~\cite{alhabash2001,paillier2020spaceground,pech2025odpll}. These effects make carrier phase recovery (CPR) a central receiver function~\cite{taylor2009phase}: the residual phase must be estimated accurately enough to preserve the decision regions of a high-order constellation as the received signal quality changes over time. Recent coherent-FSO studies consequently combine atmospheric mitigation with digital carrier synchronization or robust data-aided processing~\cite{paillier2020spaceground,johst2024dataaided,panasiewicz2022opl}.

Data-aided (DA) and non-data-aided (NDA) carrier estimators provide complementary ways to obtain that phase reference. A DA receiver inserts known pilot symbols, allowing the receiver to form direct phase observations and fit the carrier evolution from regularly spaced references~\cite{gavert2022pilot}. Data-aided processing has demonstrated stable coherent reception at very low signal-to-noise ratios (SNRs) in free-space experiments~\cite{johst2024dataaided}. An NDA phase statistic instead removes the modulation phase from the received data. The classical Viterbi--Viterbi construction raises phase-shift-keyed samples to the rotational-symmetry order before averaging their phases~\cite{viterbi1983}; the same principle extends to amplitude-phase-shift-keyed constellations by using the number of phase points on each ring~\cite{du2025nda,ma2025mapphase}, building on joint maximum-likelihood phase estimation theory developed for sinusoidal carriers under Wiener phase noise~\cite{wang2022jointml,du2021jointml,liu2025tdml}. NDA therefore retains every transmitted position for payload and can aggregate the full observation window, whereas DA obtains explicit references at a prescribed pilot density. The reported NDA bit-error evaluation resolves the resulting rotational ambiguity per processing window using the transmitted bit labels, while the selector itself uses only pre-recovery received-power statistics. Under block-varying Gamma--Gamma irradiance, the effective SNR seen by successive processing windows changes with the atmospheric gain. The two reference strategies can consequently favor different operating regions, creating a direct motivation for receiver-side selection across these regions, an approach also explored in pilot-aided adaptive receiver designs~\cite{wang2025pilotmlsd}.

This paper develops a blockwise DA/NDA selector for a coherent satellite--ground FSO link using \((8,8)\) 16-ary amplitude-phase shift keying (16APSK). The selector applies one fixed two-stage decision procedure to every 256-sample digital-signal-processing window. First, a coefficient-of-variation boundary obtained from offline AWGN received-power characterization screens the raw received-power samples. For the remaining windows, a received-power-based channel proxy is converted to an effective SNR and compared with a fixed 13~dB decision level. Both stages use one configuration fixed before evaluation and shared across all turbulence regimes. This information flow uses only quantities available before either carrier-recovery branch produces decoded bits. The resulting study makes three contributions. First, it states reproducible signal, block, pilot, and total-energy definitions that distinguish the 100-symbol atmospheric channel block from the 256-sample receiver window. Second, it develops an interpretable two-stage rule that combines window-level power variation with a blind effective-SNR estimate and evaluates the resulting selector without regime-specific retuning. Third, it quantifies the complementary operating regions of the fixed DA and NDA estimators on both data-symbol-SNR and total-energy coordinates and evaluates selector output through a common-payload bit-error-rate (BER) ratio comparison. In the predefined high-SNR work region, the data-symbol-SNR gaps are approximately 1.3, 1.2, and 1.9~dB for strong downlink, moderate uplink, and strong uplink turbulence, respectively. Including the 1.249~dB DA pilot-energy term gives corresponding total-energy-equivalent NDA advantages of approximately 2.5, 2.4, and 3.1~dB. Separately, over the evaluated downlink points, the selector provides up to about 1.3~dB common-payload BER-ratio reduction relative to fixed NDA.

The remainder of the paper is organized as follows. Section~II defines the signal impairments, atmospheric and receiver block structure, pilot pattern, and SNR convention. Section~III describes the DA and NDA estimators, the two-stage selection statistics, and the fixed decision parameters. Section~IV presents the simulation setup, the complementary DA/NDA operating regions, and the selector evaluation. Section~V concludes the paper.
